
\chapter{Typing and Extensions}

We covered the untyped \(\lambda\)-calculus, highlighting its strengths and limitations.
This section introduces the notion of typing, which restricts its expressiveness while improving the tractability of studying its properties.
The idea of typing arises from the usual definition of a function, understood as a mapping between elements of sets.
For example, consider a function  
\( f : A \to B \) defined by \( a \mapsto f(a) \), where we recognize \(A\) as the domain and \(B\) as the codomain.
Notice how the mapping is meaningful only within this context--otherwise \( f \) would lack sense.
Hence, the important of restricting what a function can take as an input as well as what it may yield as an output.

In the context of the \(\lambda\)-calculus, types serve as a means of restriction, similar ho how we use logical predicates to define sets.
Just as \( f : \mathbf{?} \to \mathbf{?} \) is refined to \( f : A \to B \) by specifying requirements on the arguments and ensuring properties of the outputs, we use types to impose requirements on the arguments of an abstraction and to ensure properties of the reduction process of that term.

Since the \(\lambda\)-calculus treats all terms as first-class citizens, attempting to impose  
restrictions directly on arguments is not effective. Instead, we define restrictions  
on all terms by pairing each term with additional information called type.  
This tells us about the nature of a given term, which may be of two kinds: a simple  
type or arrow type \(\to\). Introducing a typing system into the \(\lambda\)-calculus splits \(\Lambda\) into two categories:
Typable and Non-typable terms. These sets are disjoint. This simple type system classifies expressions such as  
\( x (\lambda y.y) \) or \( Y \) as untypable, while \( \lambda x.x \) has type  
\(\alpha \to \alpha\). Typing ensures that terms are meaningful: only normalizing terms receive a type, while non-normalizing or ill-defined terms do not.

Out of the many type systems developed for the $\la$-calculus, we focus on a small but expressive core. We begin with minimal typing, which corresponds to the simply typed $\la$-calculus with only function types. This system captures the essence of typed functional computation: terms are either variables, abstractions, or applications, and their types are built solely from base types and arrows.

We then extend this core with product types and sum types, which allow richer data structures. Product types model tuples or pairs, supporting construction and projection, while sum types model alternatives, supporting case analysis. These extensions move the type system closer to what is found in practical programming languages, while preserving the foundational simplicity of the $\la$-calculus.

\section{\centering Minimal Typing}
The untyped version of the \lcalc \ relies on abstraction and application. Minimal typing captures this simply by using base types and arrow types.

\begin{definition}The set of minimal simple types $\Tset$ is defined inductively as follows:
\begin{itemize}
\item (Type variable) If $\alpha \in \Bset$, then $\alpha \in \Tset$.
\item (Arrow type) If $\sigma, \tau \in \Tset$, then $(\sigma \to \tau) \in \Tset$.
\end{itemize}
where $\Bset$ \ is a countably infinite set of base types.
In abstract syntax:
\[
\Tset ::= \Bset \mid \Tset \to \Tset
\]
\end{definition}
\begin{note}
  These minimal types are the most basic types one can apply to the \lcalc, this is with respect to extended types, that use products and sums to add expressiveness.
\end{note}
\begin{notation}
  In general, greek letters $\alpha, \beta, \gamma,\, \dots$ are standard to denote type names.
\end{notation}
\begin{notation}
  The convention is that $\to$ is right associative, that is: $\alpha \to \gamma \to \alpha$ is the same as $(\alpha \to (\gamma \to \alpha))$.
\end{notation}
\begin{example}
Examples of simple types include:
\(
\qquad \gamma, \quad (\beta \to \gamma), \quad ((\gamma \to \alpha) \to (\alpha \to (\beta \to \gamma)))
\)
\end{example}
%\begin{definition} Notation for the Simply Typed \lcalc. We explicitly leave \texttt{<type>} unspecified.
%\begin{align*}
%    \texttt{<term>} &\;\texttt{::=}\; \texttt{<variable>} \\
%                    &\;\texttt{|}\; \texttt{"$\lambda$"}\ \texttt{<variable>}\ \texttt{":"}\ \texttt{<type>}\ \texttt{"."}\ \texttt{<term>} \\
%                    &\;\texttt{|}\; \texttt{<term>}\ \texttt{<term>} \\
%                    &\;\texttt{|}\; \texttt{"("}\ \texttt{<term>}\ \texttt{")"} \\
%  \texttt{<variable>} &\;\texttt{::}\in\; \Vset
%\end{align*}
%\end{definition}

\begin{definition} A context $\Gamma$ is a finite mapping from variables to simple types, defined inductively by:
\[
\Gamma ::= \cdot \mid \Gamma, x : \sigma
\]
where
\begin{itemize}
\item $\cdot$ is the empty context
\item $x$ is a variable
\item $\sigma \in \Tset$ is a simple type
\item $\Gamma, x : \sigma$ extends $\Gamma$ by mapping $x$ to $\sigma$
\item \( \mathrm{dom}(\Gamma) = \{\, x \mid x : \sigma \in \Gamma \,\}\)
\item \( \mathrm{rng}(\Gamma) = \{\, \sigma \mid x : \sigma \in \Gamma \,\} \)
\end{itemize}
A context is often represented as a finite sequence:
\(
  \Gamma = x_1 : \sigma_1,\, x_2 : \sigma_2,\, \dots,\, x_n : \sigma_n
\)
with all $x_i$ distinct.
\end{definition}

With this in mind we distinguish two alternative methods for typing. Typing à la Church, where types are annotated alongside $\lambda$-terms. And Typing a la Curry, where typing judgements are provided externally through some inference algorithm. Church's style of typing adds syntax on top of the already defined \lcalc. It can be said that while Church's typing introduces rules to generate well-typed terms, Curry's approach infers types from untyped expressions.
\begin{itemize}
\item Church-style:
  \(
  \quad
    \lambda x : \sigma . x
    \vdash \lambda x : \sigma . x : \sigma \to \sigma
  \).
\item Curry-style:
  \(
    \quad
    \lambda x . x
    \vdash \lambda x . x : \sigma \to \sigma
  \).
\end{itemize}
Note how Curry's style allows us to assign different valid types to terms. For example, $(\alpha \to \beta) \to (\alpha \to \beta)$ could be a valid type assignment for $\la x . \, x$, do not confuse this with  polymorphism.

\begin{remark}
The distinction between Church-style and Curry-style typing has practical consequences for implementation. 
In Church-style, every abstraction carries its type annotation explicitly, so type checking is straightforward. 
In Curry-style, terms are written without annotations, and type information must be reconstructed by solving constraints generated during typing. 
This requires the introduction of type variables and unification algorithms in order to infer consistent typings.
Hence Curry-style systems motivate the development of type inference mechanisms, while Church-style systems rely on explicit annotations.
\end{remark}


\subsubsection{\centering Typing à la Church}

Church-style typing embeds type annotations directly in terms, making typing syntax-directed and explicit. We define the following rules that tell whether a term is well-typed.


\begin{definition} We define the following typing relation:
  \label{def:church-types}
  \[
    \inferrule*[right=Var]{x \not\in \mathrm{dom}(\Gamma)}
    {\Gamma, x : \tau \vdash x : \tau}
    \quad
    \inferrule*[right=Abs]
    {x \not\in \mathrm{dom}(\Gamma) \quad \Gamma, x : \sigma \vdash M : \tau}
    {\Gamma \vdash \lambda x . M : \sigma \to \tau}
  \]
  \[
    \inferrule*[right=App]
    {\Gamma \vdash M : \sigma \to \tau \\ \Gamma \vdash N : \sigma}
    {\Gamma \vdash M \, N : \tau}
  \]
\end{definition}


  
\begin{example}
In this example, we verify that 
\(
  (\lambda f.\,\lambda x.\,f\,x)(\lambda y.\,y)
\)
is a valid term of type \(\alpha \to \alpha\) using Church-style typing.
\begin{footnotesize}
\[
  \inferrule*[right=App]
  {
    \inferrule*[right=Abs]
    {
      \inferrule*[right=Abs]
      {
        \inferrule*[right=App]
        {
          \inferrule*[right=Var]
          { }
          { \Gamma, f : \alpha \to \alpha, x : \alpha \vdash f : \alpha \to \alpha }
          \quad
          \inferrule*[right=Var]
          { }
          { \Gamma, f : \alpha \to \alpha, x : \alpha \vdash x : \alpha }
        }
        { \Gamma, f : \alpha \to \alpha, x : \alpha \vdash f\,x : \alpha }
      }
      { \Gamma, f : \alpha \to \alpha \vdash \lambda x.\,f\,x : \alpha \to \alpha }
    }
    { \Gamma \vdash \lambda f.\,\lambda x.\,f\,x : (\alpha \to \alpha) \to (\alpha \to \alpha) }
    \quad
    \inferrule*[right=Abs]
    {
      \inferrule*[right=Var]
      { }
      { \Gamma, y : \alpha \vdash y : \alpha }
    }
    { \Gamma \vdash \lambda y.\,y : \alpha \to \alpha }
  }
  {
    \vdash (\lambda f.\,\lambda x.\,f\,x)(\lambda y.\,y) : \alpha \to \alpha
  }
\]
\end{footnotesize}
\end{example}


Church-style annotations allow type inference to be straightforward because every abstraction carries its parameter type explicitly, this comes at the cost of requiring annotations on all binders, making the language more verbose.

\begin{example} The following, is an implementation of types   \ref{def:church-types}. Then we implement type annotations on abstractions:
\begin{minted}{Haskell}
  data Type
      = Base Natural
      | Arr Type Type
  data TypedTerm
      = Var String
      | Abs String Type TypedTerm
      | App TypedTerm TypedTerm
\end{minted}
Then we add the notion of context as a map from variable names to types.
\begin{minted}{Haskell}
  type Context = [(String, Type)]
\end{minted}
We define \texttt{infer}:
\begin{minted}{Haskell}
  infer :: Context -> TypedTerm -> Maybe Type
  infer ctx (Var x) =
    lookup x ctx
  infer ctx (Abs x t body) =
    case infer ((x, t) : ctx) body of
      Just bodyType -> Just (Arr t bodyType)
      Nothing -> Nothing
  infer ctx (App t1 t2) =
    case infer ctx t1 of
      Just (Arr argType resType) ->
        case infer ctx t2 of
          Just t2Type | t2Type == argType -> Just resType
          _ -> Nothing
      _ -> Nothing
\end{minted}

\begin{minted}{shell-session}
  ghci> infer [] (Abs "x" (Base 0) (Var "x"))
  Just 0->0
\end{minted}
Our Church style inference is smart enough to provide us with statements such as:
\begin{minted}{shell-session}
  ghci> infer [("y", (Base 0))] (App (Abs "x" (Base 0) (Var "x")) (Var "y"))
  Just 0
\end{minted}
provided that the context carries sufficient information and the expressions are well formed and properly annotated.

\end{example}

\begin{definition}
If $\Gamma \vdash M : \sigma$, then we say that $M$ has type $\sigma$ in $\Gamma$. A term $M$ is typable if there exist $\Gamma$ and $\sigma$ such that $\Gamma \vdash M : \sigma$.
\end{definition}

\subsubsection{\centering Typing à la Curry}

This typing flavor is probably the more elegant version of the two alternatives presented. It provides complete separation between terms and types, meaning there is no need to extend the syntax of \lcalc \ to include type annotations in abstractions. Importantly, in the monomorphic case, if a term is indeed typable, there exists an algorithmic procedure to infer a type for it without requiring explicit type annotations. This approach can be further improved in more advanced systems, sometimes referred to as ``Algorithm W'', which systematically finds types for typable terms in the polymorphic setting. This idea leads to the Hindley-Milner type system, which allows type inference for terms with polymorphic types.

We implement Typing à la Curry following \cite{NazarethNipkow1996AlgorithmW}. Algorithm \ref{algo:W-mono} presents a faithful synthesis of their monomorphic version of Algorithm W. We provide a similar alternative that is fully interoperable with the rest of our implementation; however, we omit the code to maintain clarity and readability.

\begin{algorithm}
\caption{Monomorphic Algorithm W}
\label{algo:W-mono}
\begin{algorithmic}[1]
  \Function{W}{expr, env, n}
  \State \textbf{match} expr \textbf{with}
  
  \State \quad \textbf{Var} $x$:
  \If{$x \in \text{dom}(env)$}
  \State \Return $\text{Ok}(\text{id}, env(x))$
  \Else
  \State \Return $\text{Fail}$
  \EndIf

  \State \quad \textbf{Abs} $e$:
  \State $(s, t, m) \gets W(e, env[x \mapsto \text{TVar}(n)], n+1)$
  \State \Return $\text{Ok}(s, \text{TVar}(n) \to t, m)$

  \State \quad \textbf{App} $e_1~e_2$:
  \State $(s_1, t_1, m_1) \gets W(e_1, env, n)$
  \State $(s_2, t_2, m_2) \gets W(e_2, s_1(env), m_1)$
  \State $u \gets \text{mgu}(s_2(t_1), t_2 \to \text{TVar}(m_2))$
  \State \Return $\text{Ok}(u \circ s_2 \circ s_1, u(\text{TVar}(m_2)), m_2+1)$
  \EndFunction
\end{algorithmic}
\end{algorithm}

Here, \textbf{Var $x$} and \textbf{Abs $e$} are self-explanatory. For \textbf{App $e_1~e_2$}, the algorithm infers types for the function and argument, unifies them using \texttt{mgu}, and returns the resulting type along with the composed substitutions. \texttt{TVar(n)} generates fresh type variables, \texttt{mgu} computes the most general unifier, and $\circ$ denotes substitution composition.


\begin{example} Using inference algorithms like the one earlier we can provide terms with suitable types without need for explicit type annotations:
  Let $\sigma, \tau, \rho$ be arbitrary types. Then:
  \begin{itemize}
  \item $\vdash \la x . x : \sigma \to \sigma$
  \item $\vdash \la x . \la y . x : \sigma \to \tau \to \sigma$
  \item $\vdash \la x . \la y . \la z . x \, z \, (y \, z) 
    : (\sigma \to \tau \to \rho) \to (\sigma \to \tau) \to \sigma \to \rho$
  \end{itemize}
  from our program:
  \begin{minted}{shell-session}
  ghci> infer [] (Abs "x" (Var "x")) 0
  Just ([],(a0 -> a0),1)
  ghci> infer [] (Abs "x" (Abs "y" (Var "x"))) 0
  Just ([],(a0 -> (a1 -> a0)),2)
\end{minted}
$\la x . \la y . \la z . x \, z \, (y \, z)$ is omitted for readability. Clearly $\texttt{a0} \text{ corresponds to } \alpha$ and $\texttt{a1} \text{ to } \tau$. Note how the typing context is always empty at first. For testing purposes, let us try and find a type for the $Y$ combinator:
\begin{minted}{shell-session}
  ghci> infer [] ycomb 0
  Nothing
\end{minted}

\end{example}

\begin{note}
  Typing style is largely a practical issue, a matter of taste even, since ultimately both approaches yield valid typings. The more we annotate, the less we have to infer, and the converse is also true. For the sake of simplicity, we therefore adopt typing à la Church, keeping type information explicit and straightforward.
\end{note}

\subsubsection{\centering Properties of Minimal Typing}
In \ref{sec:normalization-and-cr} we analyzed the untyped \lcalc \ as a computation model. We also provided evidence that it is missing some desirable properties, for instance, not all terms reached a normal form, as seen in \ref{sec:normalization}, some computations can run forever. This makes reasoning about programs more difficult and motivates the study of typed variants. In the following, we provide a few of the most relevant properties of this minimally typed variant, that wile restricting expressiveness, add safety and decidability. Proofs are omitted, refer to \cite{nederpelt2014:type-theory} and \cite{Barendregt1992}.

\begin{definition}
\label{def:one-step-beta}
\leavevmode
The one-step $\beta$-reduction relation $\curly$ on terms of $\Lambda^\to$ is the smallest relation satisfying:
\begin{itemize}
\item (Basis)
  \(
    (\lambda x : \sigma.\, M)\,N \;\curly\; M[x := N]
  \)
  \item (Compatibility) If \(M \curly M'\) then
  \(
    M\,N \curly M'\,N,
    N\,M \curly N\,M',
    \lambda x : \sigma.\,M \curly \lambda x : \sigma.\,M'
  \)
\end{itemize}
and similarly for all subterms.
\end{definition}

\begin{theorem}
\label{thm:church-rosser}
The Church–Rosser property holds for minimal typing.
\end{theorem}

\begin{lemma}
\label{lem:subject-reduction}
If \(\Gamma \vdash L : \rho\) and \(L \to_\beta L'\), then
\(
  \Gamma \vdash L' : \rho.
\)
\end{lemma}

\begin{theorem}
\label{thm:strong-normalisation}
Every well-types term \( M \) in minimally typed \lcalc \ is strongly normalising.
\end{theorem}

\begin{remark}
  Proofs to these statements are common-ground among most of the \lcalc \ resources cited, see \cite{nederpelt2014:type-theory}. A proof to the normalization theorem is given \ref{proof:weak-normalization} for the product extension.
\end{remark}
\section{\centering Extended Typing--Products}
As discussed, minimal typing is the simplest form of typing. It follows from this statement that it can be extended. Here we introduce products that can be thought of as records or structs:

\begin{definition} Extended types $ \Tset^\times $ is defined inductively as:
\[
  \Tset^\times ::= \Bset \mid \Tset^\times \to \Tset^\times \mid \Tset^\times \times \Tset^\times \mid 1
\]
where $\Bset$ is the set of base types.
\end{definition}


\begin{definition} The set of raw terms \( \La^\times \) is defined by:
  \[
    \La^\times ::= x \mid \lambda x : \Tset.\ \La^\times \mid \La^\times \, \La^\times \mid \langle \La^\times, \La^\times \rangle \mid \pi_1 \La^\times \mid \pi_2 \La^\times \mid *
  \]
  Here:
  \begin{itemize}
      \item \( x \in \Vset \), a countable set of variables
  \item \( \lambda x : \Tset.\ \La^\times \) is function abstraction with type annotation
  \item \( \langle \La^\times, \La^\times \rangle \) is the term constructor for product values
  \item \( \pi_1 \) and \( \pi_2 \) are projections on product terms
  \item \( * \) is the unique term of type \( 1 \)
  \end{itemize}
\end{definition}

\begin{definition}
Given $t \in \La^\times$, typing judgments are of the form \( \Gamma \vdash t : \Tset \), where \( \Gamma \) is a context \( x_1 : T_1, \dots, x_n : T_n \). The typing rules for \( \La^\times \) are:

\[
  \inferrule*[right=app]
  {\Gamma \vdash m : a \to b \\ \Gamma \vdash n : a}
  {\Gamma \vdash m\,n : b}
  \quad
  \inferrule*[right=Abs]
  {\Gamma, x : A \vdash M : B}
  {\Gamma \vdash \lambda x : A.\ M : A \to B}
\]
\[
  \inferrule*[right=Pair]
  {\Gamma \vdash M : A \\ \Gamma \vdash N : B}
  {\Gamma \vdash \langle M, N \rangle : A \times B}
  \quad
  \inferrule*[right=Proj$_1$]
  {\Gamma \vdash M : A \times B}
  {\Gamma \vdash \pi_1 M : A}
  \quad
  \inferrule*[right=Proj$_2$]
  {\Gamma \vdash M : A \times B}
  {\Gamma \vdash \pi_2 M : B}
\]
\[
  \inferrule*[right=Unit]
  { }
  {\Gamma \vdash * : 1}
  \quad
  \inferrule*[right=Var]
  { }
  {\Gamma, x : A \vdash x : A}
\]

\end{definition}

\subsubsection{\centering Reduction Rules}
The following presents the core operational semantics for the extended types with products and unit. We naturally extend $\beta$-reduction for products and projections. We the do the same for $\eta$ reduction. It is important to
 close this root rules under a context to form a full reduction relation.

\subsubsection{\centering $\beta$-reduction}

The core $\beta$-reduction rules for simply typed $\lambda$-calculus with products and unit are:

\[
  (\lambda x : A.\ M)\ N \;\curly\; M[x:=N]
\]
\[
  \pi_1 \langle M, N \rangle \;\curly\; M
  \qquad
  \pi_2 \langle M, N \rangle \;\curly\; N
\]
\[
  \langle \pi_1 M,\, \pi_2 M\rangle = M
\]
\[
  \lambda x : U.\, (M\,x) = M \quad (x \not\in \mathrm{FV}(M))
\]

\subsubsection{\centering $\eta$-laws}

The $\eta$-laws capture extensionality principles, identifying terms with their canonical expansions:

\[
  \lambda x : A.\ M\,x \;\equiv_\eta\; M \quad (x \notin FV(M))
\]
\[
  \langle \pi_1 M, \pi_2 M \rangle \;\equiv_\eta\; M
\]
\[
  M \;\equiv_\eta\; * \quad (M : 1)
\]

\noindent
\textbf{Closure:}  
Unlike $\beta$, $\eta$ is usually treated as an equational rule rather than a reduction. It is automatically closed under contexts, since definitional equality is taken as a congruence.

\subsubsection{\centering $\beta\eta$-conversion}

The combined theory, called $\beta\eta$-conversion, is the smallest congruence $\equiv_{\beta\eta}$ satisfying:

\[
  M \curly N \;\;\Rightarrow\;\; M \equiv_{\beta\eta} N
  \qquad
  M \equiv_\eta N \;\;\Rightarrow\;\; M \equiv_{\beta\eta} N
\]

That is, $\beta\eta$-conversion identifies terms up to both computation and extensionality, closed under all term formers.

\section{\centering Church Rosser Revisited}
One important theorem that fails for $\beta\eta$-reduction in the simply-typed 
$\la$-calculus with product types is the Church–Rosser property. The culprit is the 
$\eta_1$-rule. For example, let $x$ be a variable of type $A \times 1$, and consider 
the term 
\[
M = \langle \pi_1 x, \pi_2 x \rangle.
\] 
By the $  \langle \pi_1 M, \pi_2 M \rangle \;\equiv_\eta\; M $ rule, $M$ reduces to $x$, whereas by the $1 \equiv_\eta$ rule, 
$M$ reduces to $\langle \pi_1 x, *\rangle$. Both resulting terms are normal forms, 
so there is no common reduct. Hence, the Church–Rosser property fails:
\[
\langle \pi_1 x, \pi_2 x \rangle 
\;\overset{\eta}{\longrightarrow}\; \langle \pi_1 x, * \rangle \qquad
\langle \pi_1 x, \pi_2 x \rangle 
\;\overset{\eta}{\longrightarrow}\; x.
\]
As stated in \cite{selinger2008:lambdanotes}, there are two natural ways to restore the Church–Rosser property:
\begin{itemize}
  \label{items:restorech}
  \item If we consider only $\beta$-reductions and omit all $\eta$-rules, 
  confluence holds. Computationally, this is sufficient: $\beta$-reductions 
  drive all actual computation, while $\eta$-reductions merely simplify terms. 
  Moreover, $\eta$-reductions never block $\beta$-redexes and always reduce 
  term size. Hence, any $\beta$-normal form can be further reduced to a 
  $\beta\eta$-normal form, though not necessarily uniquely.

  \item Alternatively, if we remove the type $1$ and the constant $*$ from the 
  calculus, the Church–Rosser property holds even for $\beta\eta$-reduction.
\end{itemize}

\section{\centering Normalization Revisited}
\label{sec:typed-normalization}

The same way we added types and products to the \lcalc, we now update our previous definition of normal form to adequately adjust to our new requirements:
\begin{definition} A term is normal if none of its subterms is of any of the following forms:
  \label{def:prod-normal}
\[
\pi_1 \langle u, v \rangle \quad
\pi_2 \langle u, v \rangle \quad
(\lambda x: u . v)\, u
\]
since otherwise it would be reducible.
\end{definition}

\begin{remark}
  Definitions for non-normalizing, weakly normalizing and strongly normalizing terms are kept the same, only now we attend to \ref{def:prod-normal} as definiton of normal form.
\end{remark}

\begin{theorem}
\label{thm:sn-stlc-products}
Every well-typed term in the simply typed $\lambda$-calculus with product types and unit type is strongly normalizing. Formally, if $\Gamma \vdash M : A$, then there is no infinite sequence of reductions starting from $M$:
\[
M \;\curly\; M_1 \;\curly\; M_2 \;\curly\; \dots
\]
\end{theorem}
\begin{proof}
  An machine checked proof for strong normalization in Agda can be found at \cite{urciuoli2020strong}.
\end{proof}

\begin{corollary}
  The simply typed \lcalc is weakly normalizing.
\end{corollary}
  
The proofs are lengthy and well known. They are covered in full detail by \cite{girard1989:proofs-and-types} chapters $4$ and $6$. The strong normalization theorem is difficult to prove. In contrast, the proof of the weak normalization theorem is not as hard and provides measure of complexity based on the number of redexes of a given degree within a term.

\subsubsection{\centering Proof Sketch of the Weak Normalization Theorem for Product Extension}
\label{proof:weak-normalization}
To prove the weak normalization theorem, we assign a numerical measure of complexity to terms based on the degrees of their redexes. By carefully analyzing how these degrees behave under substitution and reduction, we construct an inductive argument showing that every term can be reduced, step by step, to a normal form
\begin{definition}
The degree $\partial(T)$ of a type $T$ is defined as:
\begin{itemize}
\item $\partial(T_i) = 1$ if $T_i$ is atomic
\item $\partial(U \times V) = \partial(U \to V) = \max(\partial(U), \partial(V)) + 1$
\end{itemize}
\end{definition}
\begin{definition}
The degree $\partial(r)$ of a redex \ref{def:prod-normal} $r$ is defined by:
\begin{itemize}
    \item $\partial(\pi_1\langle u, v \rangle) = \partial(\pi_2\langle u, v \rangle) = \partial(U \times V)$, where $U \times V$ is the type of $\langle u, v \rangle$
    \item $\partial((\lambda x. v)\, u) = \partial(U \to V)$, where $U \to V$ is the type of $(\lambda x. v)$
\end{itemize}
\end{definition}
\begin{definition}
\label{def:term-degree}
The degree $d(t)$ of a term $t$ is the supremum of the degrees of the redexes it contains. For any redex $r$, we have
\[
\partial(r) \le d(r)
\]
since \( d(r) = \sup\bigl(\partial(r),\, \{\partial(r') \mid r' \subset r\}\bigr) \)
\end{definition}

\paragraph{Key Lemmas}{We cover the necessary lemmas to give the reader an intuition of how the proof develops. We omit the proofs since they are already covered in detail by \cite{girard1989:proofs-and-types}
}
\begin{lemma} Degree and Substitution:
\label{lem:degree-substitution}
If $x$ is of type $U$, then
\[
  d(t[x := u]) \leq \max(d(t),\, d(u),\, \partial(U))
\]
The effect substitution has on degree does not blow up and has a limited upper bound. Not only that, we also know the upper bound.
\end{lemma}
\begin{lemma} Degree and Conversion:
\label{lem:degree-conversion}
If $t \to u$, then
\(
d(u) \le d(t)
\).
Reduction only decreases or maintains degree.
\end{lemma}
\begin{lemma} Conversion of Maximal Degree:
\label{lem:max-degree-conversion}
Let $r$ be a redex of maximal degree $n$ in $t$, and suppose all redexes strictly contained in $r$ have degree less than $n$. If $u$ is obtained by converting $r$ to $c$, then $u$ has strictly fewer redexes of degree $n$. Meaning that the a term of maximal degree reduces to a term of smaller degree.
\end{lemma}
Putting the lemmas together, we use the notion of measure to finalize the proof sketch:
\begin{definition} The Measure of a term
\label{def:term-measure}
$t$, is defined as
\[
  \mu(t) = (n, m)
\]
where: \( n = d(t) \text{ and  } m = \text{number of redexes of degree } n. \)
\end{definition}
\begin{theorem}
\label{thm:weak-normalisation} By iteratively applying Lemma~\ref{lem:max-degree-conversion} to reduce the measure $\mu(t)$ in lexicographic order, every term $t$ eventually reaches a normal form. Hence, weak normalization holds.
\end{theorem}
\section{\centering Extended Typing--Sums}
Sum types can be thought of as disjoint unions, they are the dual to product types, see \ref{def:coproduct}. Whereas a product type $A \times B$ contains both an $A$ and a $B$, a sum type $A + B$ contains either an $A$ or a $B$, together with a tag indicating which side it came from. This makes sum types a natural way to represent alternatives or branching data. As done in \cite{selinger2008:lambdanotes} we only present the extension in the effort avoid being repetitive.
\begin{definition} 
Extended types $ \Tset^+ $ is defined inductively as:
\[
  \Tset^+ ::= \Bset \mid \Tset^+ \to \Tset^+ \mid \Tset^+ + \Tset^+ \mid 0
\]
where $\Bset$ is the set of base types.
\end{definition}

\begin{definition}
The set of raw terms \( \La^+ \) is defined by:
\[
  \La^+ ::= x 
    \mid \lambda x : \Tset.\ \La^+ 
    \mid \La^+\,\La^+ 
    \mid \iota_1 \La^+ 
    \mid \iota_2 \La^+ 
    \mid \mathsf{case}\; \La^+\; \mathsf{of}\;\iota_1 x \Rightarrow \La^+ \mid \iota_2 y \Rightarrow \La^+
    \mid \square_A\,\La^+
\]
Here:
\begin{itemize}
  \item \( x \in \Vset \), a countable set of variables
  \item \( \iota_1 M \) injects a term of type $A$ into the left summand of $A+B$
  \item \( \iota_2 M \) injects a term of type $B$ into the right summand of $A+B$
  \item \( \mathsf{case}\;M\;\mathsf{of}\;\iota_1 x \Rightarrow N \mid \iota_2 y \Rightarrow P \) performs case analysis on a sum
  \item \( \square_A M \) eliminates an impossible term of type $0$, producing a term of type $A$
\end{itemize}
\end{definition}
\begin{definition}
Given $t \in \La^+$, typing judgments are of the form \( \Gamma \vdash t : \Tset \), where \( \Gamma \) is a context \( x_1 : T_1, \dots, x_n : T_n \). The typing rules for \( \La^+ \) are:
\[
  \inferrule*[right=Inl]
  { \Gamma \vdash M : A }
  { \Gamma \vdash \iota_1 M : A + B }
  \quad
  \inferrule*[right=Inr]
  { \Gamma \vdash M : B }
  { \Gamma \vdash \iota_2 M : A + B }
\]
\[
  \inferrule*[right=Case]
  { \Gamma \vdash M : A + B \\ \Gamma, x:A \vdash N : C \\ \Gamma, y:B \vdash P : C }
  { \Gamma \vdash \mathsf{case}\;M\;\mathsf{of}\;\iota_1 x \Rightarrow N \mid \iota_2 y \Rightarrow P : C }
  \quad
  \inferrule*[right=Abort]
  { \Gamma \vdash M : 0 }
  { \Gamma \vdash \square_A M : A }
\]

\end{definition}
\subsubsection{\centering Reduction Rules}
The operational semantics for sum types naturally extend $\beta$-reduction with case distinction and the empty type.
\subsubsection{\centering $\beta$-reduction}
The core $\beta$-reduction rules for simply typed $\lambda$-calculus with sums and empty type are:
\[
  \mathsf{case}\;(\iota_1 M)\;\mathsf{of}\;\iota_1 x \Rightarrow N \mid \iota_2 y \Rightarrow P \;\curly\; N[x := M]
\]
\[
  \mathsf{case}\;(\iota_2 M)\;\mathsf{of}\;\iota_1 x \Rightarrow N \mid \iota_2 y \Rightarrow P \;\curly\; P[y := M]
\]
\[
  \square_A M \;\curly\; \text{(no rule, $M:0$ cannot reduce)}
\]
\subsubsection{\centering $\eta$-laws}
The $\eta$-laws for sums capture the universal property of coproducts:

\[
  \mathsf{case}\;M\;\mathsf{of}\;\iota_1 x \Rightarrow \iota_1 x \mid \iota_2 y \Rightarrow \iota_2 y \;\equiv_\eta\; M
\]
\[
  \square_A M \;\equiv_\eta\; M \quad (M:0)
\]
As with products, $\eta$ is treated as an equational principle, closed under all contexts.
\subsubsection{\centering $\beta\eta$-conversion}
The combined theory, called $\beta\eta$-conversion, is the smallest congruence $\equiv_{\beta\eta}$ satisfying:
\[
  M \curly N \;\;\Rightarrow\;\; M \equiv_{\beta\eta} N
  \qquad
  M \equiv_\eta N \;\;\Rightarrow\;\; M \equiv_{\beta\eta} N
\]
Thus, $\beta\eta$-conversion identifies terms up to computation and extensionality, closed under sums and the empty type.
\begin{remark}
  The term constructors for sums have the following interpretation:
  \begin{itemize}
  \item $\iota_1$ and $\iota_2$ are the canonical injections, marking whether a value belongs to the left or right summand of $A + B$.
  \item The $\mathsf{case}$ construct corresponds to case analysis on a coproduct, analogous to pattern matching in functional programming.
  \item The term $\mathsf{\square}_A$ expresses the unique morphism $0 \to A$, capturing the fact that from an impossible value one may derive a term of any type.
  \end{itemize}
\end{remark}
\section{\centering Extended Typing--Products and Sums}
The systems $\La^\times$ and $\La^+$ may be combined into a single calculus
$\La^{\times+}$ with corresponding type system $\Tset^{\times+}$. In this setting
the formation rules for types, raw terms, typing rules, and reduction rules are
the union of those previously introduced for sums and for products. 
Thus,
\begin{definition}
  Extended types with products and sums, denoted \( \Tset^{\times+} \), are defined inductively as:
  \[
    \Tset^{\times+} ::= \Bset 
    \mid \Tset^{\times+} \to \Tset^{\times+} 
    \mid \Tset^{\times+} \times \Tset^{\times+} 
    \mid \Tset^{\times+} + \Tset^{\times+} 
    \mid 1 
    \mid 0
  \]
  where \( \Bset \) is the set of base types.
\end{definition}
\begin{definition}
  The set of raw terms with products and sums, denoted \( \La^{\times+} \), is given by
  \[
    \La^{\times+} ::= \La^\times \cup \La^+ .
  \]
\end{definition}
\begin{remark}
  By uniting sums and products we obtain the full simply typed $\lambda$-calculus
  with finite types. Every finite type can be expressed from $1$, $0$, sums, and
  products. The typing and operational rules are exactly those already given,
  taken together without modification.
\end{remark}
\begin{remark}
  Again, $\eta$-reduction breaks Church-Rosser. Fixes for this were given in \ref{items:restorech}.
\end{remark}
\subsubsection{\centering Kinds of problem in Type Theory}
These are the usual problems type theorists face. A programmer should be at least familiar with typability, and even more so with Type Checking, the one responsible for the daily headaches of bug finding in typed languages:
\begin{itemize}
\item \textbf{Typability}:
  Determine whether a term is legal. Formally, given a term, find a context and type such that
  \( ? \;\vdash\; \texttt{term} \, : \; ? \) holds. If the term is not well-typed, identify where the derivation fails. Type Assignment is a variant where the context is given, so the task is only to find the type: \( \Gamma \;\vdash\; \texttt{term} \, : \; ? \)
\item \textbf{Type Checking}: 
  Given a context, term, and type, check whether \( \Gamma \;\vdash\; \texttt{term} : \texttt{type} \)
  holds. The task is purely to verify that the term has the specified type in the given context.
\item \textbf{Inhabitation}: 
  Given a context and a type, determine whether there exists a term with that type: \( \Gamma \;\vdash\; ? : \texttt{type}  \)
\end{itemize}
\begin{remark}
  In programming languages and theorem provers, these type problems have concrete meanings. Well-typedness asks whether a term or program is valid under the typing rules, ensuring correctness of proofs or programs. Type Checking verifies that a given term has a specified type, allowing compilers or proof assistants to confirm correctness. Term Finding constructs a term of a given type, corresponding to automatic proof generation in theorem provers or program synthesis in programming languages.
\end{remark}
\begin{note}
  In our specific case, the implementations given show an increasing level of type inference. The simplest kind is capable of type checking and relies on annotations for every term, the second implements Church--style typing, and the third implements proper inference on monomorphic types given no annotations. The point of this note is to relate type checking and typability through annotations vs. inference, meaning the more type annotations we add, the more typability reduces to checking.
\end{note}
\paragraph{References \& Resources}{Notions that are ubiquitous to Type Theory have been outsourced from either \cite{nederpelt2014:type-theory}, \cite{selinger2008:lambdanotes} or \cite{sorensen2006:curry-howard}. The subsection on normalization is a synthesis of what can be found in \cite{girard1989:proofs-and-types}. Other texts that might have influenced the development of this section have been cited accordingly. In regard to implementation details, similar implementations to ours have been examined, namely \cite{NazarethNipkow1996AlgorithmW}, seen before, and \cite{Ahmed2012TypeInference}. The latter covers formal verification of ``Algorithm W'' and type inference in the \textbf{ML} Language.
}





